% ═════════════════════════════════════════════════════════════════════════════
\chapter{Gli algoritmi quantistici}
\label{chap:algoritmi}
% ═════════════════════════════════════════════════════════════════════════════

% ─────────────────────────────────────────────────────────────────────────────
\section{Che cos'è un algoritmo quantistico}
\label{sec:alg-intro}
% ─────────────────────────────────────────────────────────────────────────────

I capitoli precedenti hanno descritto le macchine: come si costruisce un qubit,
con quali tecnologie, dove si trovano i computer quantistici italiani. Una
macchina però, da sola, non risolve nulla. Serve sapere \textit{cosa farle
fare}, ed è questo il compito degli algoritmi.

Un algoritmo quantistico segue quasi sempre lo stesso schema in tre passi:

\begin{enumerate}
  \item \textbf{Preparazione.} I qubit partono da $\ket{0}$ e vengono messi in
    sovrapposizione, di solito con una porta di Hadamard su ciascuno. Con $n$
    qubit si ottiene una combinazione di tutti i $2^n$ valori possibili.
  \item \textbf{Evoluzione.} Una sequenza di porte modifica le ampiezze di
    questi valori. L'idea chiave è l'\textbf{interferenza}: le ampiezze delle
    risposte sbagliate si cancellano tra loro, quelle della risposta giusta si
    rinforzano.
  \item \textbf{Misura.} Si leggono i qubit. Il risultato è uno solo, scelto a
    caso secondo le probabilità finali: se l'algoritmo è ben costruito, la
    risposta giusta è la più probabile.
\end{enumerate}

\begin{notabox}
  Un'idea sbagliata molto diffusa è che il computer quantistico ``provi tutte
  le soluzioni in parallelo'' e restituisca la migliore. Non è così: alla
  misura si ottiene un solo risultato, e se le ampiezze fossero tutte uguali
  sarebbe un risultato a caso. Tutto il lavoro di un algoritmo quantistico sta
  nel far \textit{interferire} le possibilità in modo che, alla fine, quella
  giusta abbia la probabilità più alta.
\end{notabox}

Il capitolo procede in tre tappe. Prima si ricordano i due algoritmi
``storici'', Grover e Shor, già citati nel Paragrafo~\ref{sec:vantaggi}. Poi si
passa agli algoritmi pensati per le macchine di oggi, piccole e rumorose. Infine
si mostra, con un esempio concreto preso da MathWorks, come un problema reale -
organizzare i giri di consegna di una flotta di furgoni - possa essere
riscritto in una forma adatta a un computer quantistico.

% ─────────────────────────────────────────────────────────────────────────────
\section{Gli algoritmi classici della computazione quantistica}
\label{sec:alg-classici}
% ─────────────────────────────────────────────────────────────────────────────

\subsection{Grover: cercare un ago nel pagliaio}
\label{subsec:alg-grover}

Si immagini un elenco di $N$ elementi non ordinato, in cui uno solo soddisfa
una certa condizione. Un computer classico non ha scorciatoie: deve
controllarli uno alla volta, e in media ne guarda $N/2$. L'algoritmo di
Grover~\cite{grover1996fast} trova l'elemento giusto con circa $\sqrt{N}$
controlli.

Il funzionamento si può descrivere a parole. Si parte da una sovrapposizione
uniforme, in cui ogni elemento ha la stessa ampiezza. Poi si ripetono due
operazioni:

\begin{itemize}
  \item l'\textbf{oracolo} riconosce l'elemento cercato e ne cambia il segno
    dell'ampiezza, senza dire quale sia;
  \item la \textbf{diffusione} riflette tutte le ampiezze attorno al loro
    valore medio: quella col segno invertito, che sta sotto la media, finisce
    molto sopra.
\end{itemize}

A ogni ripetizione l'ampiezza dell'elemento cercato cresce un po'. Dopo circa
$\frac{\pi}{4}\sqrt{N}$ ripetizioni è vicina a 1, e la misura restituisce
quasi certamente la risposta giusta. La Figura~\ref{fig:grover} mostra lo
schema del circuito.

\begin{figure}[ht]
  \centering
  \begin{quantikz}
    \lstick{$\ket{0}$} & \gate{H} & \gate[3]{\text{Oracolo}}
      & \gate[3]{\text{Diffusione}} & \meter{} \\
    \lstick{$\ket{0}$} & \gate{H} & & & \meter{} \\
    \lstick{$\ket{0}$} & \gate{H} & & & \meter{}
  \end{quantikz}
  \caption{Schema dell'algoritmo di Grover su tre qubit. Il blocco
    oracolo + diffusione viene ripetuto circa $\frac{\pi}{4}\sqrt{N}$ volte
    prima della misura.}
  \label{fig:grover}
\end{figure}

Il guadagno è \textbf{quadratico}, non esponenziale: per un elenco di un milione
di voci si passa da circa mezzo milione di controlli a circa ottocento. È un
miglioramento notevole, ma non trasforma un problema impossibile in uno facile.

\subsection{Shor: fattorizzare i numeri}
\label{subsec:alg-shor}

L'algoritmo di Shor~\cite{shor1994algorithms} è il motivo per cui i computer
quantistici sono diventati famosi anche fuori dai laboratori. Scompone un
numero intero nei suoi fattori primi in tempo polinomiale, mentre i migliori
metodi classici richiedono un tempo che cresce quasi esponenzialmente con il
numero di cifre.

L'idea è trasformare la fattorizzazione in un altro problema: trovare il
\textbf{periodo} di una funzione, cioè ogni quanto essa si ripete. Trovare
periodi è esattamente ciò in cui la \textit{trasformata di Fourier quantistica}
eccelle. Una volta noto il periodo, i fattori si ricavano con semplici calcoli
classici.

La conseguenza pratica è nota: la sicurezza di RSA si basa sul fatto che
fattorizzare numeri di migliaia di bit è troppo difficile. Un computer
quantistico abbastanza grande la renderebbe inutile. ``Abbastanza grande'',
però, significa milioni di qubit fisici con correzione degli errori: molto
oltre le macchine descritte nei Capitoli~\ref{chap:implementazioni}
e~\ref{chap:italia}, che arrivano a qualche centinaio.

% ─────────────────────────────────────────────────────────────────────────────
\section{Gli algoritmi per le macchine di oggi}
\label{sec:alg-nisq}
% ─────────────────────────────────────────────────────────────────────────────

Grover e Shor richiedono circuiti lunghi e qubit quasi perfetti. Le macchine
attuali, chiamate \textbf{NISQ} (\textit{Noisy Intermediate-Scale Quantum},
cioè ``di media scala e rumorose'')~\cite{preskill2018quantum}, hanno invece
pochi qubit e perdono la coerenza in fretta. Per sfruttarle si è sviluppata una
famiglia di algoritmi diversi, detti \textbf{ibridi} o \textbf{variazionali}.

\subsection{L'idea ibrida}
\label{subsec:alg-ibrido}

Il principio è dividere il lavoro tra due computer:

\begin{itemize}
  \item il \textbf{computer quantistico} esegue un circuito corto, che
    dipende da alcuni parametri regolabili (per esempio gli angoli di
    rotazione delle porte), e misura il risultato;
  \item il \textbf{computer classico} guarda il risultato, valuta quanto è
    buono e sceglie nuovi parametri per il giro successivo.
\end{itemize}

Il ciclo si ripete finché il risultato smette di migliorare. È lo stesso
meccanismo con cui si addestra una rete neurale: si aggiustano dei parametri
un passo alla volta. Il vantaggio è che il circuito resta corto, e quindi
sopravvive al rumore. È anche il motivo per cui un sistema come SOL viene
installato accanto al supercomputer Leonardo (Paragrafo~\ref{sec:it-ospitate}):
le due macchine devono parlarsi continuamente.

\subsection{VQE e QAOA}
\label{subsec:alg-vqe-qaoa}

I due algoritmi variazionali più noti sono:

\begin{itemize}
  \item il \textbf{VQE} (\textit{Variational Quantum
    Eigensolver})~\cite{peruzzo2014vqe}, che cerca lo stato di energia minima
    di una molecola. Serve in chimica e nella scienza dei materiali, ed è
    l'applicazione più vicina all'idea originale di Feynman di usare un sistema
    quantistico per simularne un altro~\cite{feynman1982simulating};
  \item il \textbf{QAOA} (\textit{Quantum Approximate Optimization
    Algorithm})~\cite{farhi2014qaoa}, che cerca soluzioni buone a problemi di
    ottimizzazione combinatoria, cioè problemi in cui bisogna scegliere la
    combinazione migliore tra moltissime possibili.
\end{itemize}

\subsection{L'annealing quantistico}
\label{subsec:alg-annealing}

Esiste poi un approccio che non usa porte logiche: l'\textbf{annealing
quantistico}~\cite{kadowaki1998annealing}. Il problema viene ``scritto''
direttamente nelle interazioni tra i qubit, in modo che la soluzione migliore
corrisponda allo stato di energia più bassa. Il sistema parte da uno stato
semplice e viene trasformato lentamente verso quello che rappresenta il
problema; se la trasformazione è abbastanza lenta, i qubit restano nello stato
di energia minima e alla fine indicano la soluzione.

Le macchine di D-Wave lavorano così, e anche la \textit{modalità analogica} di
SOL descritta nel Paragrafo~\ref{sec:it-ospitate} segue un'idea simile: i qubit
evolvono con continuità invece di eseguire una sequenza di porte.

\begin{notabox}
  QAOA e annealing hanno un punto in comune importante: entrambi vogliono il
  problema scritto nella stessa forma, detta \textbf{QUBO}. Il prossimo
  paragrafo spiega cos'è, e il successivo mostra come ci si arriva partendo da
  un problema concreto.
\end{notabox}

% ─────────────────────────────────────────────────────────────────────────────
\section{La forma QUBO}
\label{sec:alg-qubo}
% ─────────────────────────────────────────────────────────────────────────────

QUBO sta per \textit{Quadratic Unconstrained Binary Optimization}. Il nome
contiene già la definizione:

\begin{itemize}
  \item \textbf{binary}: le variabili valgono solo 0 o 1, come i bit;
  \item \textbf{quadratic}: il costo da minimizzare contiene al massimo
    prodotti di due variabili;
  \item \textbf{unconstrained}: non ci sono vincoli separati. Tutte le regole
    del problema devono essere già dentro la formula del costo.
\end{itemize}

In formule, dato un vettore di variabili binarie $x = (x_1, \dots, x_n)$, si
cerca il valore che rende più piccolo
\begin{equation}
  f(x) = \sum_{i} Q_{ii}\, x_i + \sum_{i<j} Q_{ij}\, x_i x_j ,
  \label{eq:qubo}
\end{equation}
dove i numeri $Q_{ij}$ descrivono il problema. Ogni variabile corrisponde a un
qubit, e ogni coefficiente $Q_{ij}$ diventa un'interazione tra i qubit $i$ e
$j$.

Resta da capire come si fanno sparire i vincoli. Il trucco si chiama
\textbf{penalità}: per ogni regola si aggiunge al costo un termine che vale zero
quando la regola è rispettata e diventa grande quando è violata. Per esempio,
se esattamente una tra le variabili $x_1, x_2, x_3$ deve valere 1, basta
aggiungere
\begin{equation}
  A\,(1 - x_1 - x_2 - x_3)^2 ,
  \label{eq:penalita}
\end{equation}
che è zero solo se una sola variabile è accesa. Se il numero $A$ è abbastanza
grande, nessuna soluzione ``furba'' che viola la regola può convenire.
Moltissimi problemi noti si possono scrivere in questo modo~\cite{lucas2014ising}.

% ─────────────────────────────────────────────────────────────────────────────
\section{Un esempio concreto: i giri di consegna}
\label{sec:alg-cvrp}
% ─────────────────────────────────────────────────────────────────────────────

Per vedere come funziona in pratica si segue un esempio pubblicato da
MathWorks nella documentazione di MATLAB~\cite{mathworks2026cvrp}, che a sua
volta riprende il metodo di Feld e colleghi~\cite{feld2019hybrid}.

\subsection{Il problema}
\label{subsec:alg-cvrp-problema}

Un'azienda ha un magazzino (il \textit{deposito}) e deve consegnare merce a un
certo numero di clienti. Ogni cliente chiede una quantità diversa, e ogni
furgone ha una \textbf{capacità} massima. Bisogna decidere quali clienti
assegnare a ciascun furgone e in che ordine visitarli, in modo che i furgoni
non siano mai sovraccarichi e che la strada percorsa in totale sia la più
corta possibile.

È il \textbf{CVRP} (\textit{Capacitated Vehicle Routing Problem}), un problema
di ogni giorno per corrieri, aziende di distribuzione e servizi di raccolta
rifiuti. Ed è difficile: il numero di modi possibili di organizzare i giri
cresce in modo esplosivo con il numero di clienti, e già con poche decine di
clienti provarli tutti è fuori discussione.

Nell'esempio di MathWorks i dati sono:
\begin{itemize}
  \item 24 clienti più il deposito, in posizioni casuali su una mappa;
  \item richieste dei clienti tra 100 e 2\,500 unità;
  \item capacità di ogni furgone pari a 6\,000 unità.
\end{itemize}

\subsection{Dividere il problema in due}
\label{subsec:alg-cvrp-divisione}

Il CVRP mescola due domande diverse: \textit{chi} va su quale furgone, e
\textit{in che ordine} visitarli. Il metodo le separa.

\paragraph{Primo passo: formare i gruppi.}
Si raggruppano i clienti vicini tra loro, con una regola semplice: si parte dal
cliente più lontano dal deposito, si aggiungono uno alla volta i clienti più
vicini al gruppo e ci si ferma quando il prossimo farebbe superare la capacità
del furgone. Poi si ricomincia con i clienti rimasti. Alla fine si fa un
piccolo ritocco, spostando un cliente in un altro gruppo se gli è più vicino e
se c'è ancora posto. Ogni gruppo corrisponde a un furgone. Questo passo è
svolto interamente da un computer classico.

\paragraph{Secondo passo: ordinare ogni gruppo.}
Dentro ogni gruppo resta da trovare il giro più corto che parte dal deposito,
tocca tutti i clienti e torna indietro. È il celebre \textbf{problema del
commesso viaggiatore} (TSP, \textit{Travelling Salesperson Problem}), ed è
questa la parte che viene affidata al solutore QUBO.

Il vantaggio della divisione è evidente: invece di un unico problema enorme
con 24 clienti se ne hanno cinque piccoli, ciascuno con una manciata di
clienti, molto più adatti alle dimensioni delle macchine quantistiche di oggi.

\subsection{Il commesso viaggiatore in forma QUBO}
\label{subsec:alg-cvrp-tsp}

Per scrivere il TSP con variabili 0/1 si usa una tabella. Se il giro ha $N$
tappe, si introduce una variabile $x_{i,t}$ per ogni coppia (luogo $i$,
posizione $t$ nel giro):
\[
  x_{i,t} =
  \begin{cases}
    1 & \text{se il luogo } i \text{ viene visitato come } t\text{-esima tappa},\\
    0 & \text{altrimenti.}
  \end{cases}
\]
Una soluzione valida è una tabella $N \times N$ con un solo 1 in ogni riga e
in ogni colonna, come nell'esempio della Tabella~\ref{tab:tsp-esempio}.

\begin{table}[ht]
  \centering
  \renewcommand{\arraystretch}{1.2}
  \caption{Una soluzione valida per un giro di quattro tappe: il deposito D,
    poi i clienti C, A e B. Ogni riga e ogni colonna contiene un solo 1.}
  \label{tab:tsp-esempio}
  \begin{tabular}{lcccc}
    \toprule
    & Tappa 1 & Tappa 2 & Tappa 3 & Tappa 4 \\
    \midrule
    Deposito D & 1 & 0 & 0 & 0 \\
    Cliente A  & 0 & 0 & 1 & 0 \\
    Cliente B  & 0 & 0 & 0 & 1 \\
    Cliente C  & 0 & 1 & 0 & 0 \\
    \bottomrule
  \end{tabular}
\end{table}

Il costo da minimizzare è formato da tre pezzi:
\begin{equation}
  \begin{split}
    f(x) ={} & \underbrace{A \sum_{i} \Bigl(1 - \sum_{t} x_{i,t}\Bigr)^{2}}_{\text{ogni luogo una volta}}
      + \underbrace{A \sum_{t} \Bigl(1 - \sum_{i} x_{i,t}\Bigr)^{2}}_{\text{ogni tappa un luogo}} \\
    & + \underbrace{\sum_{i \neq j} d_{ij} \sum_{t} x_{i,t}\, x_{j,t+1}}_{\text{lunghezza del giro}} .
  \end{split}
  \label{eq:tsp-qubo}
\end{equation}

I primi due termini sono penalità dello stesso tipo
dell'Equazione~\eqref{eq:penalita}: valgono zero solo se la tabella ha un 1 per
riga e un 1 per colonna. Il terzo termine somma la distanza $d_{ij}$ tra due
luoghi ogni volta che il luogo $j$ viene visitato subito dopo il luogo $i$
(dopo l'ultima tappa si torna alla prima, perché il furgone rientra al
deposito). Il peso $A$ va scelto più grande delle distanze in gioco: nell'esempio
MathWorks si parte dalla distanza massima tra due luoghi, così che violare una
regola costi sempre più di qualunque percorso.

Il prezzo di questa scrittura è il numero di variabili: servono $N^2$ qubit
per un giro di $N$ tappe. Un gruppo con cinque luoghi richiede 25 qubit, uno
con dieci ne richiede 100. Ecco perché il primo passo, che spezza il problema
in gruppi piccoli, è indispensabile.

\subsection{Risolvere e verificare}
\label{subsec:alg-cvrp-soluzione}

In MATLAB il procedimento richiede poche righe: si costruisce la matrice $Q$,
la si passa alla funzione \texttt{qubo} e si chiama \texttt{solve}.

\begin{lstlisting}[language=Matlab, basicstyle=\small\ttfamily, frame=single,
  caption={Costruzione e soluzione di un QUBO in MATLAB (schema semplificato).},
  label={lst:matlab-qubo}]
qprob = qubo(QN, c, d);   % matrice quadratica, termine lineare, costante
result = solve(qprob);    % per default: ricerca tabu, classica
x = result.BestX;         % vettore 0/1 con la soluzione
\end{lstlisting}

Dopo ogni soluzione si controlla che la tabella rispetti davvero le regole
(un 1 per riga e per colonna); se non le rispetta si ripete il calcolo, fino a
dieci tentativi. Il risultato dell'esempio sono \textbf{cinque giri}, con
carichi di 5\,000, 5\,800, 5\,400, 4\,600 e 5\,700 unità: tutti sotto il limite
di 6\,000.

\begin{notabox}
  Un dettaglio va detto con onestà: nell'esempio MathWorks il QUBO
  \textbf{non viene risolto su un computer quantistico}, ma con una
  \textit{ricerca tabu}, un metodo classico che esplora le soluzioni vicine
  evitando di tornare su quelle già viste. Il valore dell'esempio non sta nel
  solutore, ma nella \textbf{formulazione}: lo stesso QUBO, senza cambiare una
  riga, può essere passato a un annealer quantistico o tradotto in un circuito
  QAOA. Ed è proprio la formulazione la parte più difficile del lavoro.
\end{notabox}

% ─────────────────────────────────────────────────────────────────────────────
\section{Cosa aspettarsi davvero}
\label{sec:alg-realismo}
% ─────────────────────────────────────────────────────────────────────────────

L'esempio dei furgoni aiuta anche a capire dove siamo oggi. La
Tabella~\ref{tab:alg-sintesi} riassume gli algoritmi visti nel capitolo.

\begin{table}[ht]
  \centering
  \renewcommand{\arraystretch}{1.3}
  \caption{Sintesi degli algoritmi quantistici presentati.}
  \label{tab:alg-sintesi}
  \begin{tabularx}{\textwidth}{>{\bfseries}lXXl}
    \toprule
    Algoritmo & A cosa serve & Guadagno & Hardware \\
    \midrule
    Grover    & Ricerca           & Quadratico                  & Futuro \\
    Shor      & Fattorizzazione   & Esponenziale                & Futuro \\
    VQE       & Chimica, materiali & Da dimostrare              & Oggi (NISQ) \\
    QAOA      & Ottimizzazione    & Da dimostrare               & Oggi (NISQ) \\
    Annealing & Ottimizzazione    & Da dimostrare               & Oggi \\
    \bottomrule
  \end{tabularx}
\end{table}

La tabella mostra una specie di paradosso. Gli algoritmi con un vantaggio
\textit{dimostrato}, Grover e Shor, richiedono macchine che ancora non
esistono. Quelli che girano sulle macchine di oggi non hanno, per ora, un
vantaggio dimostrato rispetto ai buoni metodi classici: per un problema come
il CVRP esistono euristiche classiche molto raffinate, e battere quelle è
difficile.

Questo non rende gli algoritmi ibridi inutili. Servono a imparare a
\textbf{formulare} i problemi in modo quantistico, a capire come si comportano
le macchine reali e a preparare il terreno per quando l'hardware sarà
migliorato. È lo stesso spirito con cui sono nate le piattaforme italiane del
Capitolo~\ref{chap:italia}: aperte a ricercatori e imprese proprio perché ci
si possa esercitare oggi su problemi veri.

Anche per il diamante il discorso è lo stesso. Con pochi qubit ben controllati
e tempi di coerenza lunghi (Paragrafo~\ref{subsec:nv-coerenza}), i processori
a centri NV sono candidati naturali per circuiti variazionali piccoli, e i
collegamenti fotonici tra centri diversi potrebbero un giorno permettere di
distribuire un problema su più nodi, così come l'esempio dei furgoni lo
distribuisce su più gruppi.

% ─────────────────────────────────────────────────────────────────────────────
\section{Conclusioni del Capitolo}
\label{sec:alg-conclusioni}
% ─────────────────────────────────────────────────────────────────────────────

Un algoritmo quantistico non prova tutte le soluzioni insieme: prepara una
sovrapposizione e la fa interferire in modo che la risposta giusta diventi la
più probabile. Grover e Shor mostrano fin dove può arrivare questa idea, con
guadagni quadratici ed esponenziali, ma chiedono macchine grandi e quasi prive
di errori.

Per le macchine di oggi si usano algoritmi ibridi come VQE e QAOA, o
l'annealing quantistico, in cui il computer quantistico lavora insieme a quello
classico. Molti di questi metodi richiedono che il problema sia scritto in
forma QUBO, cioè come una funzione di variabili 0/1 in cui le regole diventano
penalità.

L'esempio dei giri di consegna ha mostrato questo percorso dall'inizio alla
fine: un problema logistico reale viene diviso in parti piccole, ogni parte
diventa un commesso viaggiatore, e ogni commesso viaggiatore diventa una
tabella di qubit. Il passo più difficile non è premere ``esegui'' su una
macchina quantistica, ma arrivare a una formulazione che quella macchina sia in
grado di capire.
